% !PRINTED 763
\weberpage{763}
%<<CONT
die nicht alle $=0$ sind, worin die Exponenten $x_1, x_2, \ldots$ von einander verschiedene algebraische Zahlen sind. Wir bezeichnen mit $\Omega$ einen \gesperrt{Normalkörper}, dem alle diese Zahlen $x_1, x_2, \ldots$ angehören, und mit $\Theta$ eine primitive Zahl dieses Körpers, ferner mit

\[
\sigma' = (\Theta,\Theta'),\ \sigma'' = (\Theta,\Theta''),\ \ldots
\]

die Substitutionen des Körpers $\Omega$. Wenn durch eine dieser Substitutionen, $\sigma'$, die Zahlen $x_1, x_2, \ldots$ in $x_1', x_2', \ldots$ übergehen, so müssen auch diese von einander verschieden sein.

Es sei jetzt $u$ eine Variable und

\weq{8}{U = A_1e^{ux_1}+A_2e^{ux_2}+\cdots}

Hierin machen wir die sämmtlichen Substitutionen $\sigma', \sigma'', \ldots$ und bezeichnen die so entstehenden Functionen mit $U', U'', \ldots$ Das Product aller dieser Functionen können wir, obwohl es sich nicht bloss um algebraische Zahlen handelt, die \gesperrt{Norm} von $U$ nennen. Dieses Product hat die Form

\weq{9}{N(U) = C_1e^{uz_1}+C_2e^{uz_2}+\cdots,}

worin die $C_1, C_2, \ldots$ gleichfalls ganze rationale Zahlen sind. Die $z_1, z_2, \ldots$ sind Zahlen des Körpers $\Omega$, die wir, wenn wir die Glieder mit gleichen Exponenten in ein einziges Glied zusammenfassen, als \gesperrt{von einander verschieden} voraussetzen dürfen. Zugleich ist wegen (7)

\weq{10}{C_1e^{z_1}+C_2e^{z_2}+\cdots = 0.}

Der Ausdruck auf der rechten Seite von (9) hat aber noch, wie seine Entstehung als Norm zeigt, die Eigenschaft, ungeändert zu bleiben, wenn irgend eine der Substitutionen $\sigma', \sigma'', \ldots$ gemacht wird. Entwickelt man aber (9) nach Potenzen von $u$, so muss diese Unveränderlichkeit von jedem Gliede dieser Reihenentwickelung gelten, also wenn $h$ ein beliebiger ganzer positiver Exponent ist, für

\[
C_1z_1^h+C_2z_2^h+\cdots;
\]

diese Summe ist aber eine Zahl des Körpers $\Omega$, und daher eine \gesperrt{rationale Zahl}. Dies können wir dahin zusammenfassen:

Bedeutet $g(z)$ irgend eine ganze Function von $z$ mit rationalen Coëfficienten, so ist

\[
C_1g(z_1)+C_2g(z_2)+\cdots
\]

eine \gesperrt{rationale Zahl}.
